\documentclass[letterpaper,10pt,conference]{ieeeconf}

\IEEEoverridecommandlockouts
\overrideIEEEmargins

\usepackage{cite}
\usepackage{amsmath,amssymb,amsfonts}
\usepackage{graphicx}
\usepackage{bm}
\usepackage{hyperref}
\usepackage{multirow}
\usepackage{booktabs}
\usepackage{xcolor}

\DeclareMathOperator{\EX}{\mathbb{E}}
\def\BibTeX{{\rm B\kern-.05em{\sc i\kern-.025em b}\kern-.08em
    T\kern-.1667em\lower.7ex\hbox{E}\kern-.125emX}}

% Upload the five frozen-result PNG/PDF files to this Overleaf directory.
% The local source paths are recorded in PAPER_FIGURE_MANIFEST.md.
\graphicspath{{Figures_R1/}}

\begin{document}

\title{Drone Interception via EDMDc-Based Model Predictive Control}

\author{Sanjay Maharjan, Darren Johnson, Animesh Bangal, Yousuf Atteia,
Tariq Hlayel, and Adrian Stein%
\thanks{S. Maharjan, D. Johnson, A. Bangal, Y. Atteia, T. Hlayel, and
A. Stein are with the Department of Mechanical and Industrial Engineering,
Louisiana State University, Baton Rouge, LA 70803, USA.
{\tt\small \{smaha12,djoh378,abanga4,yattei1,thlaye1,astein\}@lsu.edu}.}}

\maketitle

\begin{abstract}
This work presents an Extended Dynamic Mode Decomposition with control
(EDMDc) framework for quadrotor trajectory tracking and moving-target
interception. A lifted predictor is identified at 100~Hz from the response of
a nonlinear 12-state quadrotor and its cascaded attitude controller to
four-dimensional outer-loop commands: collective thrust and desired
roll, pitch, and yaw. The resulting model is embedded in a constrained
receding-horizon controller and compared with a tuned reactive cascaded PID
and a yaw-scheduled hover-linear MPC on paired, previously unseen scenarios.
Across 40 aggressive tracking cases, EDMDc-MPC obtains 0.097~m mean position
RMSE, compared with 0.774~m for PID, and avoids the severe error tail observed
for the simple linear MPC. Its advantage over linear MPC is regime dependent:
the linear controller remains highly accurate on some mild trajectories and
EDMDc has lower error in 10 of 40 paired cases. In 40 moving-target trials,
both predictive controllers satisfy a distance--relative-speed dwell
criterion in all cases; EDMDc reduces mean qualified-capture time from
3.628~s to 3.483~s. Mean EDMDc optimization time is 3.995~ms in tracking;
the 99th percentile remains below the 10~ms period, while the largest sampled
solve is 14.329~ms. The results support EDMDc as a
computationally practical nonlinear predictive model, while also identifying
the operating regimes in which a local linear model remains sufficient.
\end{abstract}

\begin{keywords}
UAV interception, Koopman operator, data-driven modeling, EDMDc, model
predictive control
\end{keywords}

\section{Introduction}

The rapid proliferation of unmanned aerial vehicles (UAVs) has introduced
significant challenges for public safety and critical-infrastructure
protection. Autonomous interception of non-cooperative drones has therefore
become an active research area, with approaches ranging from visual servoing
\cite{yang_autonomous_2020,yang_high-speed_2025} and competition-driven
capture systems \cite{garcia_autonomous_2020} to reinforcement learning
\cite{darwish_drone_2023} and proportional-navigation methods
\cite{pliska_towards_2024}. These approaches demonstrate the difficulty of
combining rapid relative motion, limited actuation, and real-time onboard
computation.

Quadrotor dynamics are nonlinear and strongly coupled. Cascaded PID remains
common because of its simplicity \cite{wang_dynamics_2016}, whereas linear
quadratic and model predictive controllers introduce prediction and explicit
constraints \cite{ahmad_simulation_2020,kouvaritakis_model_2016}. Nonlinear
MPC can model aggressive flight accurately, but its online nonlinear program
may be difficult to solve at high rate
\cite{merabti_nonlinear_2015,magni_nonlinear_2009,schwenzer_review_2021}.

Koopman-based identification offers an intermediate representation: nonlinear
state evolution is approximated by linear dynamics in a lifted observable
space. Dynamic Mode Decomposition and EDMD approximate this operator from
snapshot data \cite{schmid_dynamic_2022,williams_datadriven_2015}; EDMDc
extends the construction to controlled systems. Prediction quality depends on
the dictionary and excitation data \cite{li_extended_2017}, and convergence
results apply under appropriate sampling assumptions
\cite{korda_convergence_2018}. Recent work has considered learned
eigenfunctions \cite{folkestad_extended_2020} and analytically constructed
observables \cite{netto_analytical_2021}.

This paper studies whether a structured EDMDc model improves closed-loop
performance when it replaces only the outer-loop prediction model of a
Pixhawk-like cascaded controller. Relative to our earlier fixed-yaw
formulation, the principal revisions are: (i) a complete 12-state model that
includes yaw and yaw rate; (ii) identification from aligned
$(x_k,u_k,x_{k+1})$ samples using the outer commands actually presented to the
inner controller; (iii) evaluation at 100~Hz on fresh paired tracking and
interception cases; and (iv) a qualified interception metric requiring both
proximity and low relative speed for a dwell interval. EDMDc is compared with
a tuned reactive PID and a deliberately simple hover-linear MPC. The goal is
not to claim universal superiority over linear control, but to characterize
where nonlinear lifted prediction changes the outcome.

\section{System Modeling}
\label{sec:system_modeling}

\subsection{Nonlinear Quadrotor}

The X-configuration quadrotor state is
\begin{align}
x=\begin{bmatrix}p_W^\top&v_W^\top&\eta^\top&\omega_B^\top\end{bmatrix}^\top
\in\mathbb{R}^{12},
\end{align}
where $p_W=[x\;y\;z]^\top$, $v_W=[v_x\;v_y\;v_z]^\top$,
$\eta=[\phi\;\theta\;\psi]^\top$, and
$\omega_B=[p\;q\;r]^\top$. With
$R(\eta)=R_z(\psi)R_y(\theta)R_x(\phi)$, the translational and rotational
dynamics are
\begin{subequations}
\begin{align}
\dot p_W &= v_W,\\
\dot v_W &= m^{-1}\!\left(R(\eta)e_3T-k_vv_W+f_{\mathrm{ext}}\right)-ge_3,\\
\dot\eta &= W(\phi,\theta)\omega_B,\\
\dot\omega_B &= I^{-1}\!\left(\tau-k_\omega\omega_B-omega_B\times I\omega_B\right).
\end{align}
\end{subequations}
The rotor allocator maps four nonnegative rotor forces to
$[T,\tau_\phi,\tau_\theta,\tau_\psi]^\top$ and applies motor feasibility
limits before the wrench reaches the plant. The requested and applied wrenches
are logged separately. The present nominal study sets $f_{\mathrm{ext}}=0$;
no external force is applied in the nominal study.

\subsection{Cascaded Command Interface}
\label{sec:cascade}

The compared outer-loop controllers do not command rotor speed or body torque
directly. They generate
\begin{align}
u_k=[T_k,\phi_k^*,\theta_k^*,\psi_k^*]^\top,
\label{eq:outer_input}
\end{align}
which is passed through the same cascaded attitude/rate controller, motor
allocator, and nonlinear plant. Thus, EDMDc identifies the response of the
combined inner-loop/plant subsystem,
\begin{align}
(x_k,u_k)\longmapsto x_{k+1}.
\end{align}
This interface mirrors offboard attitude control on a flight controller while
retaining the full yaw dynamics. Simulations initialize at the first reference
yaw and cap path yaw-rate commands at 0.8~rad/s to avoid artificial yaw steps.

\subsection{EDMDc Identification}
\label{sec:edmdc}

Let $z_k=\Phi(x_k)$ be a vector of observables. EDMDc fits
\begin{align}
z_{k+1}=Az_k+Bv(x_k,u_k)
\label{eq:edmdc}
\end{align}
by regularized least squares on aligned snapshot triples. Physical states are
the first 12 lifted coordinates, so $\hat x_k=C_z z_k$.

The selected 42-dimensional observable dictionary contains the 12 physical
states, sine and cosine of all Euler angles, attitude--rate and
velocity--attitude products, quadratic velocity/rate/attitude terms,
body-frame velocity, the body thrust direction, and a constant. This
physics-structured dictionary represents the gravity--attitude coupling,
coordinate rotation, angular coupling, and affine trim without a neural
network. States and inputs are standardized using statistics computed only
from the training split.

The raw input has four components, but the regression uses the
16-dimensional feature vector
\begin{align}
v(x,u)=\big[&u^\top,;(T R(\eta^*)e_3)^\top,
\sin\eta^{*\top},\cos\eta^{*\top},
(\eta^*-\eta)_{\mathrm{wrap}}^\top\big]^\top.
\label{eq:input_lift}
\end{align}
Consequently, $A\in\mathbb{R}^{42\times42}$ and
$B\in\mathbb{R}^{42\times16}$. The yaw error is wrapped to
$[-\pi,\pi)$. This construction resolves the ambiguity between the four
physical commands optimized by MPC and the 16 features presented to the
identified predictor.

\section{Control Formulation}
\label{sec:control_formulation}

\subsection{EDMDc-MPC}

At each 10~ms update, EDMDc-MPC minimizes
\begin{align}
J=&\sum_{i=0}^{N-1}\|C_zz_i-x_i^{\mathrm{ref}}\|_Q^2
+\sum_{i=0}^{N_c-1}\|u_i-u_{\mathrm{trim}}\|_R^2 \nonumber\\
&+\sum_{i=0}^{N_c-1}\|u_i-u_{i-1}\|_{R_d}^2
+\|C_zz_N-x_N^{\mathrm{ref}}\|_{Q_f}^2,
\label{eq:mpc_cost}
\end{align}
subject to \eqref{eq:edmdc}, collective-thrust limits, and
$|\phi^*|,|\theta^*|\le45^\circ$. Because the input features in
\eqref{eq:input_lift} depend on state and command, one sequential convex
linearization is solved at each control update using OSQP and a warm start.
The selected tracking controller uses $N=50$ (0.5~s), $N_c=1$, and a terminal
weight twice the stage-state weight. Its physical state-weight diagonal before
standardization is
\begin{align}
q=&[180,180,220,16,16,20,0.2,0.2,4,0.05,0.05,1],
\end{align}
obtained by applying the selected position and velocity multipliers to the
base weights. The selected input multiplier is 0.1, and the input-rate
penalty multiplier is 300. The command trust region is 12~N in thrust and
0.45~rad in desired roll and pitch. Controller parameters are chosen on
validation trajectories before the fresh confirmation experiments.

\subsection{Baselines and Information Contract}

The first baseline is a tuned cascaded PID that uses only the current
position/velocity reference and contains no trajectory feedforward. It is
therefore a reactive baseline. The second is a yaw-scheduled hover-linear MPC:
the translational model uses the local small-angle gravity coupling and is
rotated by the current yaw. It receives the same future kinematic reference as
EDMDc-MPC. Both predictive controllers use the same outer-command interface,
plant, actuator limits, update rate, and paired references, but each is tuned
on a separate validation set. Task-specific gains are used for tracking and
interception. We therefore compare two predictive models under matched
preview, while the PID comparison quantifies the benefit of the complete
predictive architecture rather than model fidelity alone.

\subsection{Qualified Interception}

For interception, the reference horizon is formed from the target's current
kinematic state. Merely crossing a capture sphere at high relative speed is
not considered a successful rendezvous. A qualified capture occurs at the
last sample of the first 20-sample window satisfying
\begin{align}
\|p_I-p_T\|_2\le0.75~\mathrm{m},\qquad
\|v_I-v_T\|_2\le3~\mathrm{m/s}.
\label{eq:capture}
\end{align}
At 100~Hz, this sampled convention corresponds to a 0.2~s dwell requirement.
An unsuccessful 20~s episode is assigned 20~s in the failure-aware mean.

\section{Simulation Protocol}
\label{sec:simulation_setup}

\subsection{Identification Data and Split}

The base dataset contains 300 closed-loop runs: 50 helix, 50 figure-eight,
50 Lissajous, 50 waypoint, 30 hover-excitation, and 70 bounded closed-loop
yaw-PRBS runs. Base runs last 100~s at 100~Hz. Sixty additional 60~s runs span
commanded tilt bands of 7.5, 15, 22.5, and 30 degrees. Mixing truncates runs to
the common 60~s duration, producing 360 runs. Training uses 342 runs and
2,051,658 aligned transitions; nine runs are reserved for validation and nine
for prediction testing. The original five validation indices
$[38,58,128,154,209]$ and test indices $[39,59,129,155,210]$ are retained,
with four high-tilt runs added to each split. The ridge penalty
$\lambda=7.5$ is selected on validation rollouts.

The PRBS excitation commands bounded yaw rate through the closed-loop inner
controller. It does not substitute desired attitude for the plant input.
Alongside $u_k$ in \eqref{eq:outer_input}, the dataset records the requested
wrench and motor-feasible applied wrench for actuator diagnostics. Every
transition retains the alignment $(x_k,u_k,x_{k+1})$.

\subsection{Prediction and Closed-Loop Tests}

The model is assessed with rolling 2~s open-loop predictions. On the five
standard held-out families, position RMSE ranges from 0.007 to 0.047~m and
velocity RMSE from 0.010 to 0.063~m/s. Two of four high-tilt held-outs show
larger errors, especially in lateral velocity and yaw rate; these cases bound
the model-validity claim rather than being discarded.

Tracking confirmation uses 40 fresh 30~s trajectories---10 each from helix,
figure-eight, Lissajous, and waypoint families---at 1.75 times the nominal
reference speed. The mean reference 95th-percentile implied tilt is
18.17 degrees. Interception confirmation uses 40 fresh 20~s targets---10 each
from straight, accelerating, helix, and weaving families. All comparisons are
paired by trajectory seed. We report mean, median, 90th percentile, and worst
position RMSE for tracking; capture rate and failure-aware qualified-capture
time are reported for interception. Paired differences use 20,000-resample
bootstrap confidence intervals.

Optimization timing is measured around the solver call on an AMD Ryzen~9
9950X workstation running Windows build 26100, Python 3.10.11, and OSQP 1.0.5.
OSQP uses warm starting with verbose output and polishing disabled; all other
solver settings retain the package defaults. The reported desktop timing is
not an embedded hard real-time guarantee.

\section{Results}
\label{sec:results}

\subsection{Open-Loop Prediction}

Figure~\ref{fig:rollout} shows a representative two-second rollout of all 12
states. Position and attitude remain close to the nonlinear simulation over
the MPC-relevant horizon; the largest displayed discrepancy occurs in pitch
rate. Across the five conventional held-out families, the worst position
RMSE is 0.047~m. High-tilt tests reach 0.161~m position RMSE and reveal larger
$v_y$ and $r$ errors, so the data do not support a claim of uniform accuracy
throughout all aggressive flight.

\begin{figure}[t]
    \centering
    \includegraphics[width=\linewidth]{3_EDMDc_Rollout.png}
    \caption{Representative 2~s open-loop EDMDc prediction on the held-out
    Lissajous trajectory. The model is initialized once and propagated using
    the recorded outer commands; it is not reset at each sample.}
    \label{fig:rollout}
\end{figure}

\subsection{Observable Capacity and Learning Curve}

Table~\ref{tab:capacity} uses the same nine-run validation and nine-run test
splits for every model. The state-only and trigonometric dictionaries
underfit, whereas the 42-term physics dictionary and 56-term extension are
nearly indistinguishable. We use the predeclared parsimony rule---select the
smaller model when validation position RMSE differs by less than 0.01~m---so
the 42-term model is retained. A 20-case control-validation check with fixed
MPC weights also produced nearly identical behavior for 42 and 56 terms; both
shared the same isolated Lissajous tail. The 56-term learning curve plateaus
on the fixed test split by 10\% of the runs, although validation error improves
from 0.0655 to 0.0505~m as the training set grows. Thus added data reduce the
selection-set error but do not support a claim that all 56 terms are needed.

\begin{table}[t]
\centering
\caption{Two-Second Rolling Position RMSE [m]}
\label{tab:capacity}
\begin{tabular}{lrrcc}
\toprule
Dictionary & $n_z$ & Runs & Validation & Test\\
\midrule
State+bias & 13 & 342 & 0.1379 & 0.2016\\
State+trig. & 19 & 342 & 0.1377 & 0.2015\\
Physics & 42 & 342 & \textbf{0.0452} & 0.0776\\
Full & 56 & 34 & 0.0655 & \textbf{0.0742}\\
Full & 56 & 86 & 0.0651 & 0.0789\\
Full & 56 & 171 & 0.0581 & 0.0772\\
Full & 56 & 342 & 0.0505 & 0.0785\\
\bottomrule
\end{tabular}
\end{table}

\subsection{Aggressive Trajectory Tracking}

\begin{table}[t]
\centering
\caption{Fresh Paired Tracking Confirmation ($n=40$)}
\label{tab:tracking}
\begin{tabular}{lrrrr}
\toprule
Controller & Mean & Median & P90 & $>1$ m\\
& \multicolumn{3}{c}{Position RMSE [m]} & \\
\midrule
Reactive PID & 0.774 & 0.799 & 1.085 & 11\\
Hover-linear MPC & 1.140 & \textbf{0.027} & 4.216 & 7\\
EDMDc-MPC & \textbf{0.097} & 0.067 & \textbf{0.310} & \textbf{0}\\
\bottomrule
\end{tabular}
\end{table}

Table~\ref{tab:tracking} and Fig.~\ref{fig:tracking_summary} show the full
distribution. EDMDc-MPC beats reactive PID in all 40 paired scenarios. Its
mean paired RMSE difference is $-0.677$~m, with a 95\% bootstrap interval of
$[-0.758,-0.597]$~m. The EDMDc mean-RMSE interval is
$[0.065,0.134]$~m, and no EDMDc case exceeds 0.51~m.

The linear result is bimodal rather than uniformly poor. It is best on the
helix family (0.010~m family-mean RMSE versus 0.028~m for EDMDc) and
figure-eight family (0.038~m versus 0.068~m), but it develops a severe tail on
the Lissajous family. EDMDc has lower paired RMSE in 10 of 40 linear
comparisons, yet its mean paired advantage is 1.044~m
(95\% bootstrap interval $[0.321,1.925]$~m) because it suppresses this tail.
Figure~\ref{fig:tracking_case} therefore presents a bounded middle-ground
example: all controllers complete the waypoint trajectory, while EDMDc
reduces RMSE from 0.894~m for linear MPC to 0.078~m. These results support a
regime-dependent benefit, not universal dominance over linear MPC.

\begin{figure}[t]
    \centering
    \includegraphics[width=\linewidth]{4_Tracking_Summary.png}
    \caption{Position-RMSE distribution, empirical CDF, family means, and tail
    statistics across 40 paired aggressive trajectories. A logarithmic scale
    exposes the hover-linear controller's bimodal error distribution.}
    \label{fig:tracking_summary}
\end{figure}

\begin{figure}[t]
    \centering
    \includegraphics[width=0.95\linewidth]{5_Waypoint_Tracking.png}
    \caption{Representative fresh waypoint confirmation (run 793008).
    Position RMSE is 0.401~m for reactive PID, 0.894~m for hover-linear MPC,
    and 0.078~m for EDMDc-MPC. The reference requires a 27.7-degree
    95th-percentile implied tilt.}
    \label{fig:tracking_case}
\end{figure}

Mean tracking solve times are 1.649~ms for hover-linear MPC and 3.995~ms for
EDMDc-MPC. All three methods have zero solver failures and zero steps at which
the allocator changes the requested wrench in the accepted tracking runs.
The worst episode-level 99th percentiles are 2.997~ms and 6.784~ms,
respectively. The largest individual samples are 7.655~ms and 14.329~ms, so
the desktop experiment does not establish a hard 10~ms bound.
The larger EDMDc cost follows from its 42-state lifted prediction and
sequential input-feature linearization.

\subsection{Moving-Target Interception}

\begin{table}[t]
\centering
\caption{Nominal Qualified Interception ($n=40$)}
\label{tab:interception}
\begin{tabular}{lrrr}
\toprule
Controller & Geometric & Qualified & Failure-aware\\
& captures & captures & mean [s]\\
\midrule
Reactive PID & 40 & 38 & 13.265\\
Hover-linear MPC & 40 & 40 & 3.628\\
EDMDc-MPC & \textbf{40} & \textbf{40} & \textbf{3.483}\\
\bottomrule
\end{tabular}
\end{table}

Both predictive controllers qualify on all 40 targets
(Table~\ref{tab:interception}). EDMDc is faster than linear MPC in 28 cases,
with a mean paired difference of $-0.1445$~s and a 95\% bootstrap interval of
$[-0.251,-0.036]$~s. This is a small timing advantage, not a capture-rate
advantage. PID often reaches the capture sphere first, but its mean relative
speed at first entry is 7.205~m/s, compared with 1.212~m/s for linear MPC and
0.938~m/s for EDMDc-MPC. Consequently, PID frequently fails the low-relative-
speed dwell requirement. Figure~\ref{fig:interception_traces} distinguishes
first entry from qualified dwell completion.

\begin{figure}[t]
    \centering
    \includegraphics[width=\linewidth]{6_Interception_Traces.png}
    \caption{Representative separation and relative-speed histories for the
    four target families. Open triangles mark first sphere entry and filled
    circles mark completion of the qualified-capture dwell.}
    \label{fig:interception_traces}
\end{figure}

\begin{figure}[t]
    \centering
    \includegraphics[width=\linewidth]{7_Interception_Summary.png}
    \caption{Qualified-capture time, capture-rate sensitivity, and relative
    speed at first sphere entry for 40 paired nominal targets. Reported times
    correspond to the end of the qualifying window.}
    \label{fig:interception_summary}
\end{figure}

In the nominal interception set, both predictive controllers have zero solver
failures and zero allocator-altered steps. Their maximum per-episode
95th-percentile solve times are 5.976~ms for EDMDc and 7.500~ms for linear MPC,
below the 10~ms update period. These statistics do not establish hard
worst-case timing on embedded hardware.

\subsection{Limitations}

The results are simulation-only. An offboard hardware implementation using
the same flight-controller-style cascaded interface is in progress, but
quantitative flight tracking has not yet been completed. The qualified-capture thresholds are task
definitions and must be validated against the physical capture mechanism.
Moreover, scripted sharp and adaptive-evasion stress tests expose an important
boundary: EDMDc does not consistently outperform the tuned linear controller,
and sustained two-second low-speed capture is not achieved by EDMDc in the
sharpest tested cases. The reported sample-wise timing is measured on the
desktop test CPU and is not a hard real-time bound on embedded hardware.
These limitations motivate hardware timing, disturbance, and evasive-target
studies rather than a claim of general interception robustness.

\section{Conclusions}
\label{sec:conclusions}

A 42-state EDMDc predictor was identified at 100~Hz for a nonlinear
quadcopter controlled through collective-thrust and desired-attitude commands.
In fresh aggressive tracking tests, EDMDc-MPC reduced mean position RMSE from
0.774~m for reactive PID to 0.097~m and avoided the severe tail of a simple
hover-linear MPC. The linear controller remained best on some mild cases,
showing that EDMDc's advantage depends on excitation of nonlinear flight
dynamics. In nominal interception, EDMDc and linear MPC both qualified on all
40 targets; EDMDc provided a modest 0.1445~s mean timing advantage and lower
first-entry relative speed. Aggressive evasive tests delimit the present
model's validity. Future
work will prioritize quantitative flight validation, measured worst-case
embedded timing, and identification data that cover sharp pursuit maneuvers.

\bibliographystyle{IEEEtran}
\bibliography{references}

\end{document}
